\documentclass[11pt,a4paper]{article}
\usepackage[margin=1in]{geometry}
\usepackage{times}
\usepackage{latexsym}
\usepackage{graphicx}
\usepackage{booktabs}
\usepackage{amsmath}
\usepackage{multirow}
\usepackage{xcolor}
\usepackage{subcaption}
\usepackage{natbib}
\usepackage{hyperref}

\title{The Minimal Sufficient Representation for Verdict Prediction:\\Exception Summaries Match Full Trace Performance at 16$\times$ Compression}

\author{Anonymous}

\begin{document}
\maketitle

%==============================================================================
\begin{abstract}
Execution traces have shown promise for improving code understanding tasks, yet prior work reports inconsistent benefits and evaluates only generative settings.
We study \emph{verdict prediction}---binary pass/fail classification of code--test pairs---and ask: \textbf{what is the minimal sufficient trace representation?}
Through systematic ablation across four code LLMs (Qwen-2.5-Coder-7B, Qwen-2.5-Coder-1.5B, DeepSeek-Coder-6.7B, CodeLlama-13B-Instruct), we find that a single exception summary (${\sim}50$ tokens) matches or exceeds all richer trace representations (${\sim}800$ tokens) in every tested condition.
On Qwen-7B, exception-only achieves 99.83\% accuracy with zero variance across three seeds, matching the dataset's theoretical ceiling imposed by two irreducible samples.
Richer VOTA traces achieve statistically equivalent performance on three of four models (gap $\leq 0.42$pp) but degrade by 14.88pp on DeepSeek, where 91.7\% of training samples suffer trace truncation.
Removing explicit \texttt{[ERROR]} format markers produces identical results across all seeds, confirming that models learn from exception \emph{content}, not formatting cues.
Loop-compression signals show high seed variance ($\sigma = 3.85$pp), disqualifying them as stable predictors.
We formalize a \emph{minimal sufficient representation} principle for verdict tasks: extract the highest-specificity diagnostic signal and discard everything else.
\end{abstract}

%==============================================================================
\section{Introduction}

Execution traces---step-by-step records of program behavior---are an intuitive source of signal for code understanding tasks.
Recent work has explored trace-augmented supervised fine-tuning (SFT) for code generation \cite{ni2024next,ding2024traced}, with mixed results: DCSH \cite{ni2024next} found ``no single trace representation consistently outperforms others,'' while TracePile \cite{gu2025tracepile} scaled trace SFT to 2.6M samples but evaluated only generative tasks.

A fundamental question remains unaddressed: what trace information actually matters, and for which tasks?
We study \emph{verdict prediction}---classifying whether a code solution passes or fails a given test case---and find a striking answer: \textbf{a compact exception summary of roughly 50 tokens is sufficient to match the best trace representations of roughly 800 tokens.}

Our investigation proceeds in four stages.
First, we establish that traces substantially help verdict prediction: on Qwen-2.5-Coder-7B, compressed VOTA traces improve accuracy from 75.13\% (no trace) to 99.41\% (+24.3pp), while full uncompressed traces reach only 81.76\%.
Second, component ablation reveals that exception events alone (99.83\%) match or exceed full VOTA, while loop-compression signals (82.55\% $\pm$ 3.85pp across three seeds) are volatile and fully subsumed.
Third, cross-model evaluation on four architecturally distinct LLMs shows that exception summaries achieve $\geq$99\% on three of four models, with the fourth (DeepSeek-Coder-6.7B) degraded primarily by training-time trace truncation rather than model capability limitations.
Fourth, format ablation demonstrates that removing explicit \texttt{[ERROR]} markers has zero effect---models learn from exception content (error types and messages), not formatting cues.

These findings establish a \emph{minimal sufficient representation} principle: for verdict prediction, the highest-specificity diagnostic signal (exception events) is both necessary and sufficient; lower-specificity behavioral patterns are redundant and may introduce noise.
We further identify \emph{representation-induced model dependency}---richer representations create model-specific failure modes (truncation, capacity limitations) that simpler representations avoid entirely.

\paragraph{Contributions.}
(1) We demonstrate that exception summaries (${\sim}50$ tokens) are the minimal sufficient representation for verdict prediction, matching the best trace formats (${\sim}800$ tokens) across multiple models.
(2) We decompose the trace compression benefit into information density (74\%) and gradient dilution (26\%) through label-only loss ablation.
(3) We provide causal evidence that training-time truncation, not model capability, is the primary driver of cross-model VOTA performance gaps.
(4) We establish that format markers contribute zero signal---exception content alone drives learning.

%==============================================================================
\section{Related Work}

\paragraph{Trace-augmented code LLMs.}
DCSH \cite{ni2024next} evaluated five trace formats for code generation via SFT, finding inconsistent benefits across models and formats.
TracePile \cite{gu2025tracepile} compiled 2.6M execution traces for continued pretraining, reporting gains on code generation benchmarks but not classification tasks.
TRACED \cite{ding2024traced} showed that execution traces improve encoder-based fault localization.
All prior trace SFT work targets generative tasks or encoder-based classification; none examines decoder-based verdict prediction or performs systematic component ablation.

\paragraph{Verdict prediction.}
Binary pass/fail prediction of code--test pairs is a foundational subtask for test generation, mutation testing, and automated program repair.
Prior approaches rely on code-only features or encoder models.
We are the first to study trace SFT for decoder-based verdict prediction, revealing that the task's binary fault-detection nature creates a distinct information hierarchy compared to generative settings.

%==============================================================================
\section{Method}

\subsection{Task: Verdict Prediction}

Given a code solution $c$, a test input $x$, and an expected output $y$, the verdict prediction task is to predict the binary label $v \in \{\text{pass}, \text{fail}\}$ indicating whether executing $c$ on $x$ produces $y$.
We frame this as a classification task via SFT: the model receives $(c, x, y)$ along with an optional trace representation $t$ and generates the verdict token.

\subsection{Trace Representations}

We evaluate six trace representations of decreasing information content:

\begin{enumerate}
\item \textbf{Full trace}: Complete line-by-line execution trace including variable states, control flow, and all runtime events. Mean length ${\sim}$3000 tokens.
\item \textbf{VOTA (Verdict-Oriented Trace Abstraction)}: Compressed trace retaining three components---loop summaries, exception events with \texttt{[ERROR]} markers, and actual/expected value comparisons. Mean length ${\sim}$800 tokens.
\item \textbf{Exception-only}: Only the exception component of VOTA---error type and message for failing cases, empty for passing cases. Mean length ${\sim}$50 tokens.
\item \textbf{Exception-only (no marker)}: Exception content without the \texttt{[ERROR]} format prefix---raw error type and message only.
\item \textbf{Loop-only}: Only the loop-compression component of VOTA.
\item \textbf{No trace}: Code and test case only, no execution information.
\end{enumerate}

Additionally, we test \textbf{label-only loss}: standard full traces but with loss computed only on the verdict token (oracle gradient masking), isolating the gradient dilution effect.

\subsection{Dataset}

We construct a verdict prediction dataset from HumanEval (164 problems) and MBPP-sanitized (427 problems), totaling 591 problems.
For each problem, we collect 10 LLM-generated solutions and execute them against test cases, yielding 5,910 labeled samples (train: 4,720; test: 1,190).
The pass/fail ratio is approximately 71.5\%/28.5\%.
Train--test input overlap is zero by construction (problem-level split).

The exception-only representation exhibits 39.45\% training-set duplication (4,720 $\to$ 2,858 unique samples) because multiple passing solutions for the same test produce identical empty-exception traces.
This is an intrinsic property of the representation's compactness, not a data quality issue---raw traces show only 1.81\% duplication.

\subsection{Models and Training}

\paragraph{Core models.} Qwen-2.5-Coder-7B-Instruct and DeepSeek-Coder-6.7B-Instruct serve as the primary evaluation models, chosen for architectural diversity: distinct tokenizers, pretraining corpora (5.5T vs.\ 2T tokens), and execution reasoning capability (CRUXEval-I: 56.5\% vs.\ 39.0\%).

\paragraph{Supplementary models.} Qwen-2.5-Coder-1.5B-Instruct and CodeLlama-13B-Instruct provide additional validation, reported in supplementary analyses.

\paragraph{Training configuration.} All models are fine-tuned with LoRA (rank 16, $\alpha = 32$) for 3 epochs using AdamW (lr = $2 \times 10^{-4}$, cosine schedule).
Maximum sequence length is 2,048 tokens.
Each core experiment is run with three random seeds (42, 123, 456) to assess robustness.

%==============================================================================
\section{Experiments}

\subsection{Trace Representations Compared}
\label{sec:main-results}

Table~\ref{tab:main-results} presents the primary comparison of trace representations on Qwen-2.5-Coder-7B.

\begin{table}[t]
\centering
\small
\begin{tabular}{lcccc}
\toprule
\textbf{Condition} & \textbf{Acc.} & \textbf{ECE} & \textbf{Brier} & \textbf{Tokens} \\
\midrule
No trace        & 75.13 & 0.1715 & 0.1972 & 0 \\
Full trace      & 81.76 & 0.0958 & ---    & ${\sim}$3000 \\
Label-only loss & 86.22 & 0.1107 & ---    & ${\sim}$3000 \\
Loop-only       & 82.55{\tiny$\pm$3.85} & 0.0831 & ---    & ${\sim}$400 \\
VOTA            & 99.41 & 0.0054 & 0.0050 & ${\sim}$800 \\
Exception-only  & 99.83{\tiny$\pm$0.00} & 0.0027 & 0.0013 & ${\sim}$50 \\
AE-only         & 99.92 & 0.0004 & 0.0008 & ${\sim}$60 \\
\bottomrule
\end{tabular}
\caption{Verdict prediction accuracy (\%), expected calibration error (ECE), and Brier score on Qwen-2.5-Coder-7B. Exception-only and AE-only report mean$\pm$std across 3 seeds; loop-only reports mean$\pm$std across 3 seeds. AE-only = assertion/expected-value component only.}
\label{tab:main-results}
\end{table}

Exception-only achieves 99.83\% with perfect seed robustness (three seeds produce identical results, including the same two error samples).
This matches the dataset's theoretical ceiling: the two irreducible errors (sample IDs 901, 905) stem from Python's \texttt{GeneratorExit}---a \texttt{BaseException} subclass triggered by generator short-circuiting in \texttt{all()}, which the tracer marks as \texttt{[ERROR]} despite being benign control flow.

Full VOTA (99.41\%) is statistically significantly lower than exception-only (McNemar $p = 0.01$), indicating that non-exception components (loop summaries, actual/expected comparisons) introduce noise rather than complementary signal.

\subsection{Information Density Decomposition}

The gap between full traces (81.76\%) and VOTA (99.41\%) is 17.65pp.
Label-only loss ablation (computing loss only on the verdict token while providing full traces as input) achieves 86.22\%, decomposing the gap:

\begin{itemize}
\item \textbf{Gradient dilution}: 4.46pp (86.22\% $-$ 81.76\%), accounting for ${\sim}$26\% of the gap. Non-verdict tokens dilute gradients during training.
\item \textbf{Information overload}: 13.19pp (99.41\% $-$ 86.22\%), accounting for ${\sim}$74\%. Even with oracle loss masking, the model cannot extract the diagnostic signal from verbose traces.
\end{itemize}

This establishes information density---not gradient efficiency---as the primary bottleneck for trace-based verdict prediction.

\subsection{Cross-Model Analysis}
\label{sec:cross-model}

Table~\ref{tab:cross-model} summarizes VOTA vs.\ exception-only across all four models.

\begin{table}[t]
\centering
\small
\begin{tabular}{lccc}
\toprule
\textbf{Model} & \textbf{Exc.-only} & \textbf{VOTA} & \textbf{$\Delta$} \\
\midrule
Qwen-7B       & 99.83 & 99.41 & +0.42$^{*}$ \\
Qwen-1.5B     & 99.83 & 99.50 & +0.33 \\
CodeLlama-13B & 99.33 & 99.08 & +0.25 \\
DeepSeek-6.7B & 99.33$^\dagger$ & 84.45 & +14.88$^{*}$ \\
\bottomrule
\end{tabular}
\caption{Exception-only vs.\ VOTA accuracy (\%) across models. $^{*}$Statistically significant (McNemar $p < 0.05$). $^\dagger$Median of 3 seeds (seed42/456 = 99.33\%; seed123 = 91.09\%).}
\label{tab:cross-model}
\end{table}

On three of four models, exception-only and VOTA are statistically indistinguishable (gap $\leq 0.42$pp).
The DeepSeek anomaly (14.88pp gap) has a specific causal explanation: 91.7\% of VOTA training samples are truncated from a mean of 2,626 tokens to ${\sim}$2,048 tokens during training.
Group analysis confirms the causal direction: test samples whose training-set traces were truncated achieve 83.04\% accuracy, while intact samples achieve 100\% (Fisher exact $p \approx 0$).
Exception-only traces (${\sim}$50 tokens) are never truncated, explaining their consistent cross-model performance.

\subsection{Qwen-1.5B: Refuting the Execution Reasoning Bottleneck}

Table~\ref{tab:qwen15b} presents results for Qwen-2.5-Coder-1.5B, a model with substantially lower execution reasoning capability (estimated CRUXEval-I ${\sim}$30--35\% vs.\ 56.5\% for Qwen-7B).

\begin{table}[t]
\centering
\small
\begin{tabular}{lccc}
\toprule
\textbf{Condition} & \textbf{Qwen-1.5B} & \textbf{Qwen-7B} & \textbf{$\Delta$} \\
\midrule
Exception-only & 99.83 & 99.83 & 0.00pp \\
VOTA           & 99.50 & 99.41 & +0.09pp \\
Full trace     & 73.95 & 81.76 & $-$7.81pp \\
No trace       & 67.31 & 75.13 & $-$7.82pp \\
\bottomrule
\end{tabular}
\caption{Qwen-1.5B vs.\ Qwen-7B verdict prediction accuracy (\%). The 1.5B model matches the 7B on exception-only and VOTA despite large gaps on no-trace and full-trace conditions.}
\label{tab:qwen15b}
\end{table}

Despite a 7.82pp gap on the no-trace baseline (reflecting lower general capability), Qwen-1.5B matches the 7B model exactly on exception-only (99.83\%) and achieves statistically equivalent VOTA performance (99.50\% vs.\ 99.41\%, $p > 0.15$).
This directly refutes the hypothesis that VOTA utilization requires strong execution reasoning: a 1.5B model with CRUXEval-I ${\sim}$30--35\% achieves ${\sim}$99.5\% VOTA accuracy, demonstrating that structured diagnostic signals bypass the execution reasoning bottleneck entirely.

The cross-model gap concentrates in low-information conditions (no trace, full trace), where model capacity determines performance.
Once a structured diagnostic signal is provided---whether exception-only or VOTA---even a small model saturates.

\subsection{Format Marker Ablation}
\label{sec:nomarker}

To test whether models rely on the \texttt{[ERROR]} format marker as a shortcut, we train exception-only models with markers removed, retaining only the raw exception content (error type and message).

\begin{table}[t]
\centering
\small
\begin{tabular}{lcccc}
\toprule
\textbf{Seed} & \textbf{Marker} & \textbf{No-marker} & \textbf{$\Delta$} & \textbf{Errors} \\
\midrule
42  & 99.83 & 99.83 & 0.00pp & 2 (same) \\
123 & 99.83 & 99.83 & 0.00pp & 2 (same) \\
456 & 99.83 & 99.83 & 0.00pp & 2 (same) \\
\midrule
Mean & 99.83 & 99.83 & 0.00$\pm$0.00 & --- \\
\bottomrule
\end{tabular}
\caption{Marker vs.\ no-marker exception-only on Qwen-7B. Six independent training runs (3 seeds $\times$ 2 conditions) produce identical results, including the same two error samples.}
\label{tab:nomarker}
\end{table}

Table~\ref{tab:nomarker} shows that all six independent training runs produce identical results: 99.83\% accuracy with the same two errors (IDs 901, 905).
The \texttt{[ERROR]} format marker contributes exactly zero bits of information conditioned on exception content ($I(\text{verdict}; \text{marker} \mid \text{exception\_content}) = 0$).
Models learn the fault-detection mapping from exception semantics---error type names and messages---not from formatting cues.

\subsection{Loop-Only Seed Variability}

Loop-compression signals, the behavioral component of VOTA, show substantial training variance.

\begin{table}[t]
\centering
\small
\begin{tabular}{lccc}
\toprule
\textbf{Seed} & \textbf{Accuracy} & \textbf{95\% CI} & \textbf{Pass/Fail F1} \\
\midrule
42  & 85.63 & [83.5, 87.6] & 0.896 / 0.762 \\
123 & 78.24 & [75.9, 80.5] & 0.838 / 0.650 \\
456 & 83.78 & [81.6, 85.9] & 0.882 / 0.738 \\
\midrule
Mean$\pm$std & 82.55$\pm$3.85 & --- & --- \\
\bottomrule
\end{tabular}
\caption{Loop-only component across 3 seeds on Qwen-7B. Standard deviation of 3.85pp contrasts with exception-only's 0.00pp.}
\label{tab:loop-seeds}
\end{table}

The 3.85pp standard deviation (Table~\ref{tab:loop-seeds}) contrasts sharply with exception-only's perfect seed stability.
Loop signals encode graded behavioral patterns (iteration counts, execution path summaries) that require more complex input--verdict mappings, making learning sensitive to LoRA initialization.
Seed 123 drops to 78.24\%, close to the full-trace baseline (81.76\%), while seeds 42 and 456 cluster near 84--86\%.
This high variance disqualifies loop signals as a stable verdict predictor: they carry variable behavioral signal, not a robust diagnostic tier.

\subsection{Supplementary Model: CodeLlama-13B}

CodeLlama-13B-Instruct provides a third architectural family for validation.

\begin{table}[t]
\centering
\small
\begin{tabular}{lcc}
\toprule
\textbf{Condition} & \textbf{CodeLlama-13B} & \textbf{Qwen-7B} \\
\midrule
Exception-only & 99.33 & 99.83 \\
VOTA           & 99.08 & 99.41 \\
No trace       & 74.45 & 75.13 \\
\bottomrule
\end{tabular}
\caption{CodeLlama-13B-Instruct results. Exception-only and VOTA are statistically indistinguishable (gap = 0.25pp, $p > 0.3$).}
\label{tab:codellama}
\end{table}

CodeLlama-13B confirms the cross-model pattern (Table~\ref{tab:codellama}): exception-only (99.33\%) and VOTA (99.08\%) are statistically equivalent ($\Delta = 0.25$pp, McNemar $p > 0.3$), while the no-trace baseline (74.45\%) is consistent with other models.
The exception-only representation generalizes across Qwen (RoPE), DeepSeek, and CodeLlama (Llama-based) architectures.

\subsection{DeepSeek Seed Robustness and Truncation Analysis}

DeepSeek-Coder-6.7B reveals two important phenomena: LoRA initialization sensitivity and training-time truncation effects.

\begin{table}[t]
\centering
\small
\begin{tabular}{lccccc}
\toprule
\textbf{Condition} & \textbf{Seed} & \textbf{Acc.} & \textbf{Fail-R} & \textbf{Pass-R} \\
\midrule
\multirow{3}{*}{Exception-only} & 42 & 99.33 & 1.000 & 0.991 \\
& 456 & 99.33 & 1.000 & 0.991 \\
& 123 & 91.09 & 0.705 & 0.993 \\
\midrule
VOTA & 42 & 84.45 & --- & --- \\
No trace & 42 & 73.36 & --- & --- \\
\bottomrule
\end{tabular}
\caption{DeepSeek-Coder-6.7B results. Exception-only seed123 is a LoRA initialization outlier (fail recall 0.705), but still exceeds VOTA (84.45\%) by 6.64pp.}
\label{tab:deepseek}
\end{table}

Two of three seeds achieve 99.33\% with identical error patterns, while seed 123 drops to 91.09\% due to LoRA initialization landing in a suboptimal basin (fail recall = 0.705 vs.\ 1.0 for other seeds).
Notably, train loss is nearly identical across seeds (0.1225 vs.\ 0.122), yet accuracy differs by 8.2pp---token-level loss masks sample-level classification differences, especially for the minority fail class (28.5\%).

Even the worst seed (91.09\%) exceeds VOTA (84.45\%) by 6.64pp, confirming exception-only's relative dominance across all initializations.
The VOTA gap (14.88pp vs.\ Qwen) is primarily driven by training-time truncation: group analysis shows truncated-trace samples achieve 83.04\% vs.\ 100\% for intact samples (Fisher exact $p \approx 0$), establishing a causal link between trace truncation and performance degradation.

%==============================================================================
\section{Discussion}

\subsection{The Minimal Sufficient Representation}

Our central finding is that exception summaries of ${\sim}$50 tokens constitute the minimal sufficient representation for verdict prediction, matching ${\sim}$800-token VOTA traces across all tested conditions.
This sufficiency holds across four architecturally distinct models (Table~\ref{tab:cross-model}), across model scales (1.5B to 13B, Table~\ref{tab:qwen15b}), and is invariant to format markers (Table~\ref{tab:nomarker}).

The mechanism is straightforward: verdict prediction is fundamentally fault detection, and exception events are the highest-specificity fault indicators available in execution traces.
Passing cases produce empty exception traces; failing cases produce structured error descriptions.
This near-binary signal is learnable with minimal model capacity (even 1.5B parameters suffice) and minimal training data (2,858 unique samples yield 99.83\%).

\subsection{Representation-Induced Model Dependency}

Richer trace representations create failure modes absent in simpler ones.
VOTA's ${\sim}$800 tokens trigger training-time truncation on DeepSeek (91.7\% of samples affected), causing a 14.88pp gap that disappears entirely with 50-token exception summaries.
Full traces (${\sim}$3,000 tokens) suffer even more severe truncation, with DeepSeek's full-trace v1 collapsing entirely due to label truncation.

This \emph{representation-induced model dependency} has a practical design implication: simpler representations are more portable across models, tokenizers, and training configurations.
The exception-only representation is short enough to avoid truncation under any reasonable \texttt{max\_length} setting, eliminating an entire class of training-time failure modes.

\subsection{Loop Signals as Variable Behavioral Indicators}

Loop-compression signals carry genuine behavioral information (mean 82.55\% vs.\ 75.13\% no-trace baseline, +7.42pp), but their high seed variance ($\sigma = 3.85$pp) and full subsumption by exception signals make them unreliable as standalone predictors.
The variance likely reflects the graded nature of loop signals: unlike the near-binary exception signal, loop patterns encode continuous behavioral features (iteration counts, path summaries) that require more complex learned mappings sensitive to initialization.

\subsection{Token-Level Loss as an Unreliable Proxy}

We observe a consistent decoupling between token-level training loss and verdict accuracy.
On DeepSeek, seeds with nearly identical train loss (0.1225 vs.\ 0.122) differ by 8.2pp in accuracy.
Similarly, VOTA eval loss increases from epoch 1 to 3 (+65\%) while accuracy improves (+6.3pp).
For verdict prediction, checkpoint selection should use task-specific metrics (accuracy, F1), not token-level loss.

\subsection{Limitations}

Our study evaluates a single task (verdict prediction) on a single dataset (HumanEval + MBPP-sanitized, 591 problems).
While cross-model validation provides evidence of generalizability, the minimal-sufficient-representation principle may not transfer to generative code tasks where richer trace information could be necessary.
The DeepSeek VOTA result is confounded by training-time truncation; the relative contributions of truncation and model capability cannot be fully disentangled.
We evaluate four models from three architectural families; broader model coverage would strengthen universality claims.

%==============================================================================
\section{Conclusion}

We establish that a single exception summary of ${\sim}$50 tokens constitutes the minimal sufficient representation for verdict prediction, matching the best trace representations (${\sim}$800 tokens) across four code LLMs spanning three architectural families and two model scales.
The finding is robust to seed variation (0.00pp std on Qwen-7B), format marker removal (zero contribution), and model scale (1.5B matches 7B).

Our results reframe VOTA-style trace compression: its value lies not in generic compression but in \emph{fault indicator isolation}---automatically extracting exception events from noisy execution traces.
The minimal sufficient representation principle---extract the highest-specificity diagnostic signal, discard everything else---provides actionable guidance for trace representation design in verdict-oriented code intelligence tasks.

\bibliographystyle{plainnat}
\begin{thebibliography}{9}
\bibitem[Ni et~al.(2024)]{ni2024next}
Ansong Ni, Miltiadis Allamanis, Arman Cohan, Yinlin Deng, Kensen Shi, Charles Sutton, and Pengcheng Yin.
\newblock Next-step hint generation for introductory programming with execution traces.
\newblock In \emph{Proceedings of EMNLP Findings}, 2024.

\bibitem[Gu et~al.(2025)]{gu2025tracepile}
Alex Gu, Wen-Ding Li, Naman Jain, Theo X.~Olausson, Celine Lee, Koushik Sen, and Armando Solar-Lezama.
\newblock TracePile: A large-scale execution trace dataset for code LLMs.
\newblock In \emph{Proceedings of NeurIPS}, 2025.

\bibitem[Ding et~al.(2024)]{ding2024traced}
Yangruibo Ding, Marcus Min, Gail Kaiser, and Baishakhi Ray.
\newblock TRACED: Execution-aware pre-training for source code.
\newblock In \emph{Proceedings of ICSE}, 2024.
\end{thebibliography}

%==============================================================================
\appendix
\section{Qwen-1.5B Detailed Results}

\begin{table}[h]
\centering
\small
\begin{tabular}{lcccc}
\toprule
\textbf{Condition} & \textbf{Acc.} & \textbf{Pass F1} & \textbf{Fail F1} & \textbf{Train Loss} \\
\midrule
Exception-only & 99.83 & 0.999 & 0.997 & --- \\
VOTA           & 99.50 & 0.997 & 0.988 & 0.065 \\
Full trace     & 73.95 & 0.808 & 0.582 & 0.058 \\
No trace       & 67.31 & --- & --- & --- \\
\bottomrule
\end{tabular}
\caption{Qwen-1.5B full results ($n = 1{,}190$). VOTA--exception gap (0.33pp) is not statistically significant.}
\label{tab:qwen15b-full}
\end{table}

The information density effect is amplified at smaller scale: the full-trace--to--VOTA gap is 25.55pp on 1.5B (vs.\ 17.65pp on 7B), consistent with smaller models being more susceptible to information overload.

\section{CodeLlama-13B Detailed Results}

\begin{table}[h]
\centering
\small
\begin{tabular}{lcccc}
\toprule
\textbf{Condition} & \textbf{Acc.} & \textbf{Pass F1} & \textbf{Fail F1} & \textbf{Train Loss} \\
\midrule
Exception-only & 99.33 & 0.995 & 0.988 & 0.110 \\
VOTA           & 99.08 & 0.994 & 0.983 & 0.055 \\
No trace       & 74.45 & 0.813 & 0.600 & 0.175 \\
\bottomrule
\end{tabular}
\caption{CodeLlama-13B-Instruct full results ($n = 1{,}190$).}
\label{tab:codellama-full}
\end{table}

\section{Calibration Details}

Exception-only achieves near-perfect calibration (ECE = 0.0027, Brier = 0.0013), indicating the model not only predicts correctly but assigns appropriate confidence.
The no-trace condition shows severe miscalibration (ECE = 0.1715, Brier = 0.1972)---63$\times$ worse ECE and 150$\times$ worse Brier than exception-only.
Trace signals, especially Tier-1 diagnostic signals, fundamentally improve both accuracy and calibration.

\section{Irreducible Error Analysis}

The two errors across all Qwen exception-only runs (IDs 901, 905) come from HumanEval/95 (\texttt{check\_dict\_case}), where Python's \texttt{all()} short-circuits a generator, triggering \texttt{GeneratorExit}.
As a \texttt{BaseException} subclass (not \texttt{Exception}), this is benign control flow, but the tracer marks it as \texttt{[ERROR]}.
The model sees seven \texttt{[ERROR] GeneratorExit} lines and reasonably predicts fail, though the code is correct.
With only two such samples in the entire dataset, the model has no training signal to learn the benign nature of \texttt{GeneratorExit}.
This confirms 99.83\% as the representation's theoretical ceiling on this dataset.

\end{document}
